Given that most of the above references focus on Coxeter systems of particular types, it is perhaps
natural to wonder whether one can systematically study Kazhdan--Lusztig cells of an arbitrary
Coxeter system $(W,S)$. The \emph{$\la$-function}, a certain function $\la: W\ra \Z_{\ge 0}$, offers
an important tool in this study.  The $\la$-function is due to Lusztig, who first defined it for
Weyl and affine Weyl groups in \cite{LusztigCellI, LusztigCellII} and then extended the definition
to general Coxeter groups in \cite{LusztigHecke}.  Lusztig showed that under a certain boundedness
condition on $(W,S)$, the $\la$-function takes a constant value on every two-sided cell of $W$ and
hence on every left or right cell. Lusztig conjectures that the boundedness condition is satisfied
by every Coxeter system (see \cite[\S13.4]{LusztigHecke}), and we shall assume that the conjecture
holds throughout the paper (as is common in the literature). Under this assumption, it makes sense
to speak of the $\la$-value of a cell, and the set $W_n:=\{w\in W:\la(w)=n\}$ is a union of
two-sided cells for all $n\in \Z_{\ge 0}$.  In principle, one may also hope to organize and study
cells by $\la$-values.  Indeed, cells of $\la$-values 0 and 1 in a general Coxeter system $(W,S)$
are well understood. For example, if we assume $(W,S)$ is {irreducible} (see
Remark~\ref{rmk:irreducible}), then both $W_0$ and $W_1$ contain a single two-sided cell, we have
$W_0=\{1\}$, and $W_1$ consists of all non-identity elements in $W$ with a unique reduced word.
Moreover, Coxeter systems where $W_1$ is finite have been classified, and the cardinality of $W_1$
is known for such systems. For more details on the cell $W_1$, see Proposition~\ref{prop:facts01},
\cite{LusztigSubregular}, \cite{Hart}, \cite[Chapter 13]{BonnafeCells} and \cite{XuSubregular}. 
